\clearpage
\SourceLeaf{454}%
\pdfbookmark[0]{Gauss an Enke}{gauss-an-enke-1849}
\renewcommand{\BookRunningTitle}{Gauss an Enke}
\begingroup\brokenpenalty=10000
\begin{center}{\large GAUSS AN ENKE.}\end{center}
\BookRule{.20\textwidth}
\begin{center}Hochverehrter Freund!\end{center}

— — Ihre freundlich mitgeteilten Bemerkungen über die Häufigkeit der Primzahlen waren für mich in mehr als einer Hinsicht interessant. Sie haben mich an meine eigenen Arbeiten zu diesem Gegenstand erinnert, deren erste Anfänge weit zurückliegen, im Jahr 1792 oder 1793, als ich mir \textsc{Lambert}s Supplemente zu den Logarithmentafeln angeschafft hatte. Noch bevor ich mich mit feineren Untersuchungen der höheren Arithmetik beschäftigt hatte, gehörte es zu meinen ersten Aufgaben, meine Aufmerksamkeit auf die abnehmende Häufigkeit der Primzahlen zu richten. Dazu zählte ich sie in den einzelnen Tausendergruppen und verzeichnete die Ergebnisse auf einem der angehefteten weißen Blätter. Ich erkannte bald, dass diese Häufigkeit trotz aller Schwankungen im Durchschnitt annähernd umgekehrt proportional zum Logarithmus sei, sodass die Anzahl aller Primzahlen unter einer gegebenen Grenze $n$ annähernd durch das Integral
\[
\int\frac{\mathrm{d}n}{\log n}
\]
ausgedrückt werde, wobei der natürliche Logarithmus gemeint ist. Als ich später die in \textsc{Vega}s Tafeln (von 1796) abgedruckte Liste bis 400031 kennenlernte, dehnte ich meine Zählung weiter aus, was dieses Verhältnis bestätigte. Große Freude bereitete mir 1811 das Erscheinen von \textsc{Chernac}s
\SourceLeaf{455}%
cribrum, und ich verwendete sehr oft einzelne freie Viertelstunden darauf, bald hier, bald dort eine Tausendergruppe abzuzählen (da ich zu einer anhaltenden Zählung der Reihe nach keine Geduld hatte); schließlich ließ ich es jedoch ganz liegen, ohne mit der Million ganz fertig zu werden. Erst später machte ich mir \textsc{Goldschmidt}s Fleiß zunutze, teils um die verbliebenen Lücken in der ersten Million zu füllen, teils um die Zählung nach \textsc{Burck\-hardt}s Tafeln fortzusetzen So sind (nun schon seit vielen Jahren) die ersten drei Millionen abgezählt und mit dem Integralwert verglichen. Ich gebe hier nur einen kleinen Auszug:

\begin{center}
\begingroup\fontsize{10}{11.5}\selectfont\setlength{\tabcolsep}{3pt}
\begin{tabular}{r|r|rr|rr}
Unter & \shortstack{gibt es\\Primzahlen} & \shortstack{Integral\\$\displaystyle\int\frac{\mathrm{d}n}{\log n}$} & Abweich. & \shortstack{Ihre\\Formel} & Abweich. \\
500000 & 41556 & $41606{,}4$ & $+50{,}4$ & $41596{,}9$ & $+40{,}9$ \\
1000000 & 78501 & $79627{,}5$ & $+126{,}5$ & $78672{,}7$ & $+171{,}7$ \\
1500000 & 114112 & $114263{,}1$ & $+151{,}1$ & $114374{,}0$ & $+264{,}0$ \\
2000000 & 148883 & $149054{,}8$ & $+171{,}8$ & $149233{,}0$ & $+350{,}0$ \\
2500000 & 183016 & $183245{,}0$ & $+229{,}0$ & $183495{,}1$ & $+479{,}1$ \\
3000000 & 216745 & $216970{,}6$ & $+225{,}6$ & $217308{,}5$ & $+563{,}5$
\end{tabular}\endgroup
\end{center}

Dass sich auch \textsc{Legendre} mit diesem Gegenstand beschäftigt hatte, war mir nicht bekannt. Ihr Brief veranlasste mich, in seiner \textit{Théorie des Nombres} nachzusehen, und ich fand in der zweiten Ausgabe einige darauf bezogene Seiten, die ich früher übersehen (oder seitdem vergessen) haben muss. \textsc{Legendre} verwendet die Formel
\[
\frac{n}{\log n-A}
\]
wobei $A$ eine Konstante sein soll, für die er $1{,}08366$ setzt. Nach einer flüchtigen Rechnung finde ich damit in den obigen Fällen die Abweichungen
\[
\begin{array}{r}
-23{,}3\\
+42{,}2\\
+68{,}1\\
+92{,}8\\
+159{,}1\\
+167{,}6
\end{array}
\]

Diese Differenzen sind noch kleiner als die mit dem Integral, scheinen aber bei zunehmendem $n$ schneller zu wachsen als diese, sodass es durchaus möglich
\SourceLeaf{456}%
wäre, dass jene bei einer viel weiteren Fortsetzung die letzteren überträfen. Um Zählung und Formel in Übereinstimmung zu bringen, müsste man jeweils statt $A=1{,}08366$ setzen
\[
\begin{array}{r}
1{,}09040\\
1{,}07682\\
1{,}07582\\
1{,}07529\\
1{,}07179\\
1{,}07297
\end{array}
\]

Es scheint, dass bei wachsendem $n$ der (Durchschnitts-)Wert von $A$ abnimmt. Ob aber der Grenzwert, wenn $n$ ins Unendliche wächst, $1$ oder eine von $1$ verschiedene Größe sein wird, darüber wage ich keine Vermutung. Ich kann nicht sagen, dass es einen Grund gibt, einen ganz einfachen Grenzwert zu erwarten; andererseits könnte der Überschuss von $A$ über $1$ durchaus eine Größe der Ordnung $\frac{1}{\log n}$ sein. Ich wäre geneigt zu glauben, dass das Differential der betreffenden Funktion einfacher sein muss als die Funktion selbst. Während ich dafür $\frac{\mathrm{d}n}{\log n}$ angenommen habe, würde \textsc{Legendre}s Formel eine Differentialfunktion voraussetzen, die etwa $\frac{\mathrm{d}n}{\log n-(A-1)}$ wäre. Ihre Formel könnte übrigens für ein sehr großes $n$ als übereinstimmend mit
\begingroup\predisplaypenalty=10000\postdisplaypenalty=10000
\[
\frac{n}{\log n-\frac{1}{2k}}
\]
betrachtet werden, wobei $k$ der Modul der \textsc{Briggs}schen Logarithmen ist, also mit \textsc{Legendre}s Formel, wenn
\[
A=\tfrac{1}{2k}=1{,}1513\ \text{gesetzt wird.}
\]
\endgroup

Schließlich will ich noch bemerken, dass ich zwischen Ihren Zählungen und meinen ein paar Unterschiede bemerkt habe.
\begin{center}
\begin{tabular}{lrlrlrlr}
Zwischen & 59000 & und & 60000 & haben Sie & 95 & ich & 94\\
 & 101000 & & 102000 & & 94 & & 93
\end{tabular}
\end{center}
Der erste Unterschied liegt vielleicht daran, dass in \textsc{Lambert}s Supplementen die Primzahl 59023 zweimal aufgeführt ist. Die Tausendergruppe von 101000—102000 wimmelt in \textsc{Lambert}s Supplementen von Fehlern; in meinem Exemplar habe ich 7 Zahlen markiert, die keine Primzahlen sind, und hingegen 2 fehlende eingefügt.
\SourceLeaf{457}%
Könnten Sie nicht den jungen \textsc{Dase} dazu veranlassen, die Primzahlen in den folgenden Millionen aus den Tafeln bei der Akademie abzuzählen, die, wie ich fürchte, das Publikum nicht besitzen soll? Für diesen Fall bemerke ich, dass die Zählung in der 2. und 3. Million auf meine Anweisung nach einem besonderen Schema vorgenommen wurde, das ich selbst auch schon bei einem Teil der ersten Million angewandt hatte. Die Zählungen für jeweils 100000 stehen auf einer (kleinen) Oktavseite in 10 Spalten, von denen sich jede auf eine Zehntausendergruppe bezieht; dazu kommt noch eine Spalte davor (links) und eine dahinter rechts. Als Beispiel folgt hier eine senkrechte Spalte mit den beiden Zusatzspalten aus dem Intervall $1000000\mathinner{.\,.\,.}1100000$ \mbox{— — —}

\Needspace{20\baselineskip}
Zur Erläuterung diene zum Beispiel die 1. senkrechte Spalte. In der Zehntausendergruppe von 1000000 bis 1010000 liegen 100 Hundertergruppen; darunter ist 1, die nur eine Primzahl enthält; gar keine mit 2 oder 3; 2 mit je 4 Primzahlen; 11 mit je 5 usw. Alle zusammen ergeben $752=1\pdot1+4\pdot2+5\pdot11+6\pdot14+\dotstwo$ Die letzte Spalte enthält die Summen aus den 10 einzelnen Spalten. Die Zahlen 14. 15. 16 in der ersten senkrechten Spalte sind hier überflüssig, da keine Hundertergruppen mit so vielen Primzahlen vorkommen; auf den folgenden Blättern kommen sie aber zur Anwendung. Zuletzt werden die 10 Seiten wieder zu 1 vereinigt und umfassen so die ganze 2. Million.

Doch es ist Zeit abzubrechen. \mbox{— — —} Mit herzlichen Wünschen für Ihr Wohlbefinden
\begin{center}Stets der Ihrige\end{center}
\noindent Göttingen, 24. Dezember 1849.\hfill C. F. \textsc{Gauss}.
\BookRule{.25\textwidth}
\par\endgroup

